\section{Hauptklausur 2020}

% \newpage
\subsection*{Aufgabe 1 (12 Punkte): Kurzfragen}
\begin{enumerate}
\item Die Fourier-Reihe $F_\infty(x)$ einer $2\pi$-periodischen Funktion $f\in R([0,2\pi])$ ist gegeben durch
\[
f(x) = \sum_{k=-\infty}^{\infty} c_k e^{ikx}.
\]
Geben Sie die Integralformel der Koeffizienten $c_k$ an.\\
% Antwort:
% \[
% c_k=\frac{1}{2\pi}\int_{0}^{2\pi} f(x)\,e^{-ikx}\,\mathrm{d}x,\qquad k\in\mathbb{Z}.
% \]

\item Die Fourier-Reihe $F_\infty(x)$ einer $2\pi$-periodischen Funktion $f\in R([0,2\pi])$ sei gegeben durch
$$
f_\infty(x) = \sum_{-\infty}^\infty c_k e^{ikx}
$$
$c_k$. Geben Sie die Parsevalsche Gleichung (Vollständigkeitsrelation) mithilfe der $c_k$ an.\\
% Antwort:
% \[
% \frac{1}{2\pi}\int_{0}^{2\pi}|f(x)|^2\,\mathrm{d}x=\sum_{k=-\infty}^{\infty}|a_k|^2,
% \]
% (sofern $f\in L^2([0,2\pi])$).

\item Seien $x$ und $y$ Elemente eines reellen oder komplexen Vektorraums, der mit einem Skalarprodukt $\langle\cdot,\cdot\rangle$ versehen sei. Geben Sie die Cauchy–Schwarz–Ungleichung an.\\
% Antwort:
% \[
% |\langle x,y\rangle|\leq\|x\|\,\|y\|,
% \]
% wobei $\|x\|=\sqrt{\langle x,x\rangle}$; Gleichheit genau dann, wenn $x$ und $y$ linear abhängig sind.

\item Sei $X$ eine beliebige nichtleere Menge. Definieren Sie den Begriff der Metrik auf $X$.\\
% Antwort: Eine Abbildung $d:X\times X\to[0,\infty)$ heißt Metrik, wenn für alle $x,y,z\in X$ gilt:
% \begin{enumerate}
% \item $d(x,y)=0\iff x=y$ (Definitheit),
% \item $d(x,y)=d(y,x)$ (Symmetrie),
% \item $d(x,z)\le d(x,y)+d(y,z)$ (Dreiecksungleichung).
% \end{enumerate}

\item Sei $x=(x_1,\dots,x_n)^{\mathrm T}\in\mathbb{R}^n$. Geben Sie die Formel der \(\ell^1\)-Norm \(\|x\|_1\) an.\\
% Antwort:
% \[
% \|x\|_p=\Big(\sum_{i=1}^n |x_i|^p\Big)^{1/p},\qquad 1\le p<\infty,
% \]
% insbesondere $\|x\|_1=\sum_{i=1}^n|x_i|$, $\|x\|_2=\sqrt{\sum_{i=1}^n x_i^2}$ und $\|x\|_\infty=\max_{1\le i\le n}|x_i|$.

\item Definieren Sie den Begriff einer offenen Teilmenge $O\subseteq\mathbb{R}^n$.\\
% Antwort: $O\subset\mathbb{R}^n$ heißt offen, wenn für jedes $x\in O$ ein $\varepsilon>0$ existiert mit $B_\varepsilon(x)\subset O$, wobei $B_\varepsilon(x)=\{y\in\mathbb{R}^n:\|y-x\|<\varepsilon\}$.

\item Sei $M\subseteq\mathbb{R}^n$. Definieren Sie den Begriff der Folgenkompaktheit von $M$.\\
% Antwort: $M$ ist folgenkompakt, falls jede Folge in $M$ eine konvergente Teilfolge besitzt, deren Grenzwert wieder in $M$ liegt.

\item Sei $A\in\mathbb{R}^{n\times n}$. Definieren Sie positive Definitheit für eine solche Matrix $A$.\\
% Antwort: $A$ heißt (symmetrisch) positiv definit, wenn $A=A^{\mathrm T}$ und für alle $x\in\mathbb{R}^n\setminus\{0\}$ gilt $x^{\mathrm T}Ax>0$.

\item Sei $D\subseteq\mathbb{R}^n$ offen und $f:D\to\mathbb{R}$ stetig differenzierbar. Formulieren Sie den Mittelwertsatz für eine solche Funktion $f$.\\
% Antwort: Für $x,y\in D$ mit der ganzen Strecke $[x,y]\subset D$ existiert ein $\xi$ auf dem Segment zwischen $x$ und $y$ so dass
% \[
% f(y)-f(x)=\nabla f(\xi)\cdot (y-x).
% \]

\item Sei $f:\mathbb{R}^2\to\mathbb{R}^2$ definiert durch $f(x,y)=(x^2y^2,\;xy^4)^{\mathrm T}$. Geben Sie die Jacobi-Matrix $J_f(x,y)$ von $f$ an.\\
% Antwort:
% \[
% Jf(x,y)=\begin{pmatrix}
% \partial_x(x^2-y^2) & \partial_y(x^2-y^2)\\[4pt]
% \partial_x(2xy) & \partial_y(2xy)
% \end{pmatrix}
% =\begin{pmatrix}
% 2x & -2y\\[4pt]
% 2y & 2x
% \end{pmatrix}.
% \]

\item Sei $y'(t)=A(t)y(t),\ t\ge t_0$ ein lineares Differentialgleichungssystem für $y:[t_0,\infty)\subset \mathbb{R} \to \mathbb{R}^n$ und $\Phi(t)$ die zugehörige (Haupt-)Fundamentalmatrix. Geben Sie das Anfangswertproblem an, das durch $\Phi(t)$ gemäß Definition gelöst wird.\\
% Antwort: Die Fundamentalmatrix $\Phi(t)$ erfüllt
% \[
% \Phi'(t)=A(t)\,\Phi(t),\qquad \Phi(t_0)=I_n.
% \]
% Für ein Anfangswertproblem $y'(t)=A(t)y(t),\;y(t_0)=y_0$ ist die Lösung gegeben durch $y(t)=\Phi(t)\,y_0$.

\item Sei $D\subseteq\mathbb{R}^n$ und $F\in C(D;\mathbb{R}^n)$. Definieren Sie den Begriff des Potentials von $F$.\\
% Antwort: Eine Funktion $\varphi\in C^1(D)$ heißt Potential (oder Stammfunktion) von $F$, falls
% \[
% \nabla\varphi(x)=F(x)\quad\text{für alle }x\in D.
% \]
% \end{enumerate}

\newpage
\subsection*{Aufgabe 2 (12 Punkte)}
(a) Man bestimme mit kurzer Begründung das Komplement \(M^{c}\), den Rand \(\partial M\), das Innere \(M^{\circ}\) und den Abschluss \(\overline{M}\) der Teilmenge \(M\subseteq\mathbb{R}^{n}\) definiert durch
\[
M:=\{x\in\mathbb{R}^{n}\mid \|x\|_{\infty}\le 1,\ x\neq 0\}.
\]
Dabei bezeichnet \(\|\cdot\|_{\infty}\) die Maximumsnorm.

Bemerkung: Man achte beim Rand auf eine saubere Begründung!

(b) Sei \(\varnothing\neq A\subseteq\mathbb{R}^{n}\). Man beweise, dass der Rand \(\partial A\) kompakt ist, falls \(A\) kompakt ist. %\([3]\)

(c) Man untersuche jeweils, ob die folgenden Funktionen \(f,g:\mathbb{R}^{2}\to\mathbb{R}\) im Ursprung \((x,y)^{T}=(0,0)^{T}\) stetig sind:

(i)
\[
f(x,y)=
\begin{cases}
\dfrac{x-y}{|x|+|y|^{2}}, & (x,y)^{T}\neq(0,0)^{T},\\[6pt]
1, & (x,y)^{T}=(0,0)^{T}.
\end{cases}
\]

(ii)
\[
g(x,y)=
\begin{cases}
\dfrac{x^{2}-y^{2}}{|x|+|y|}, & (x,y)^{T}\neq(0,0)^{T},\\[6pt]
0, & (x,y)^{T}=(0,0)^{T}.
\end{cases}
\]

\newpage
\subsection*{Aufgabe 3 (12 Punkte)}
(a) Sei $f:\mathbb{R}^2\to\mathbb{R}$ eine beliebig oft stetig partiell differenzierbare Funktion, definiert durch
\[
f(x,y)=84\,y\,e^{x^2+y}.
\]
(i) Man berechne die Richtungsableitung $\partial_v f(0,0)$ von $f$ an der Stelle $(x,y)^\top=(0,0)^\top$ in die Richtung $v=\bigl(\tfrac{\sqrt{3}}{2},\tfrac{1}{2}\bigr)^\top\in\mathbb{R}^2$.

(ii) Man begründe, ob die Hesse-Matrix $H_f(x,y)$ für alle $(x,y)^\top\in\mathbb{R}^2$ symmetrisch ist, ohne diese zu berechnen.

(b) Man bestimme alle existierenden Extrema der Funktion $f:\mathbb{R}^2\to\mathbb{R}$ gegeben durch
\[
f(x,y)=(x-4)^2+(y+2)^2
\]
unter der Nebenbedingung
\[
g(x,y)=x^2+y^2-5=0
\]
mittels Lagrange-Multiplikatoren und gebe die entsprechenden Lagrange-Multiplikatoren an.

Bemerkung: Die Extrema müssen nicht klassifiziert werden.

(c) Sei $U\subseteq\mathbb{R}^n$ offen und $f:U\to\mathbb{R}$ eine stetig differenzierbare Funktion. Sei $x_0\in U$ und $f(x_0)=:c$. Betrachte $M\subseteq U$ zudem die Niveaumenge definiert durch
\[
M:=\{x\in U\mid f(x)=c\}.
\]
Sei ferner $\varepsilon>0$ und $\varphi:(-\varepsilon,\varepsilon) \subset \mathbb{R} \to\mathbb{R}^n$ eine beliebige stetig differenzierbare Abbildung mit
\[
\varphi(0)=x_0\quad\text{und}\quad \varphi(t)\in M\ \forall t\in(-\varepsilon,\varepsilon).
\]
Man zeige, dass dann
\[
\langle\varphi'(0),\nabla f(x_0)\rangle=0,
\]
wobei $\langle\cdot,\cdot\rangle$ das euklidische Skalarprodukt, $\varphi'$ die erste Ableitung von $\varphi$ und $\nabla f$ den Gradienten von $f$ bezeichnet.

\newpage
\subsection*{Aufgabe 4 (12 Punkte)}
(a) Man zeige, dass die Gleichung
\[
 x^{3}+2xy-4xz+2y-1=0
\]
in einer Umgebung \(U\subseteq\mathbb{R}^{2}\) von \((x_{0},y_{0})^{T}=(1,1)^{T}\) durch eine eindeutig stetige und in \((x_{0},y_{0})^{T}=(1,1)^{T}\) differenzierbare Funktion \(\varphi:U\to\mathbb{R}\) mit \(\varphi(1,1)=1\) nach \(z\) aufgelöst werden kann und berechne den Gradienten \(\nabla\varphi(1,1)\).

(b) Die Gleichung
\[
 e^{2x}-x=t
\]
kann in einer Umgebung \(U\subseteq\mathbb{R}\) von \(t_{0}=1\) eindeutig nach \(x\) durch eine zweimal stetig differenzierbare Abbildung \(\varphi:U\to\mathbb{R}\) mit \(\varphi(1)=0\) aufgelöst werden. Dabei gilt für die erste Ableitung, \(\varphi'(1)=1\). Man berechne die zweite Ableitung \(\varphi''(1)\).

(c) Sei \(\Omega:=\{(x,y,z)^{T}\in\mathbb{R}^{3}\mid (x,y)^{T}\neq(0,0)^{T}\}\) und \(F\) ein Vektorfeld auf \(\Omega\) definiert durch
\[
 F(x,y,z):=\left(-\frac{y}{x^{2}+y^{2}},\ \frac{x}{x^{2}+y^{2}},\ 1\right)^{T}.
\]
Weiter sei \(\gamma:[0,2\pi]\to\Omega\) eine Kurve definiert durch
\[
 \gamma(t):=(\cos t,\ \sin t,\ \sin t)^{T}.
\]
Man berechne das vektorielle Kurvenintegral \(\displaystyle\int_{\gamma}F\).

\newpage
\subsection*{Aufgabe 5 (12 Punkte)}
\begin{enumerate}
\item[(a)] Man forme die Differentialgleichung 3. Ordnung für \(y:[a,b]\subset\mathbb{R}\to\mathbb{R}\), \(a<b\), gegeben durch
\[
\frac{d^{3}}{dt^{3}}y(t)=3\frac{d^{2}}{dt^{2}}y(t)-\Bigl(\frac{d}{dt}y(t)\Bigr)^{2}+42\,y(t)+1,\qquad t\in[a,b],
\]
in ein äquivalentes Differentialgleichungssystem 1.\ Ordnung um.

\item[(b)] Es sei das Anfangswertproblem für \(y:[0,\infty)\subset\mathbb{R}\to[0,\infty)\subset\mathbb{R}\) definiert durch
\[
\frac{d}{dt}y(t)=|\sqrt{y(t)}|,\qquad t\in[0,\infty),
\]
\[
y(0)=y_{0},
\]
mit \(y_{0}\ge 0\) gegeben.
\begin{enumerate}
\item[(i)] Man bestimme unter welcher Bedingung an \(c\in\mathbb{R}\) der Ansatz
\[
y(t)=\tfrac{1}{4}(t+c)^{2}
\]
die Differentialgleichung des gegebenen Anfangswertproblems löst.

\item[(ii)] Man bestimme \(c\) in Abhängigkeit von \(y_{0}\), sodass \(y(t)=\tfrac{1}{4}(t+c)^{2}\) das gegebene Anfangswertproblem löst.

\item[(iii)] Eine weitere Lösung des Anfangswertproblems mit \(y_{0}=0\) ist \(\tilde{y}(t)=0\) für alle \(t\in[0,\infty)\). Warum widerspricht die Existenz zweier unterschiedlicher Lösungen des Anfangswertproblems mit \(y_{0}=0\) nicht dem Eindeutigkeitssatz aus der Vorlesung? Man begründe die Antwort.
\end{enumerate}

\item[(c)] Man berechne das Taylor-Polynom \(T_{3}(y,t,t_{0}=0)\) mit Entwicklungspunkt \(t_{0}=0\) der 3.\ Ordnung der beliebig oft stetig differenzierbaren Lösung \(y:[0,\infty)\to\mathbb{R}\) des Anfangswertproblems
\[
\frac{d}{dt}y(t)=\cos\bigl(y(t)\bigr),\qquad t\in[0,\infty),
\]
\[
y(0)=0.
\]
Bemerkung: Man muss dazu die Lösung \(y\) nicht tatsächlich berechnen!
\end{enumerate}

